\documentclass[12pt,a4paper]{article}
\usepackage[utf8]{inputenc}
\usepackage[spanish]{babel}
\usepackage[margin=2.5cm]{geometry}
\usepackage{graphicx}
\usepackage{hyperref}
\usepackage{titlesec}
\usepackage{enumitem}
\usepackage{fancyhdr}
\usepackage{lastpage}
\usepackage{xcolor}
\usepackage{tabularx}
\usepackage{booktabs}
\usepackage{parskip}

% Configuración de hiperenlaces
\hypersetup{
    colorlinks=true,
    linkcolor=blue,
    filecolor=magenta,      
    urlcolor=cyan,
}

% Configuración de encabezados y pies de página
\pagestyle{fancy}
\fancyhf{}
\rhead{Memoria de Prácticas Externas}
\lhead{Alejandro Vialard Santana}
\rfoot{Página \thepage\ de \pageref{LastPage}}

% Configuración de títulos
\titleformat{\section}{\Large\bfseries}{\thesection.}{1em}{}
\titleformat{\subsection}{\large\bfseries}{\thesubsection.}{1em}{}
\titleformat{\subsubsection}{\normalsize\bfseries}{\thesubsubsection.}{1em}{}

\begin{document}

%========================================
% PORTADA
%========================================
\begin{titlepage}
    \centering
    \vspace*{2cm}
    
    {\Huge\bfseries Memoria de Prácticas Externas\par}
    \vspace{1cm}
    {\Large Prácticas Curriculares\par}
    
    \vspace{2cm}
    
    \begin{figure}[h]
        \centering
        % \includegraphics[width=0.4\textwidth]{logo_uex.png}
    \end{figure}
    
    \vspace{2cm}
    
    {\LARGE\bfseries Alejandro Vialard Santana\par}
    \vspace{0.5cm}
    {\Large Grado en Ingeniería Informática en Ingeniería del Software\par}
    
    \vspace{2cm}
    
    {\large Escuela Politécnica de Cáceres\par}
    {\large Universidad de Extremadura\par}
    
    \vspace{2cm}
    
    {\large\bfseries Fecha de entrega: 27 de enero de 2026\par}
    
    \vfill
\end{titlepage}

%========================================
% ÍNDICE
%========================================
\tableofcontents
\newpage

%========================================
% 1. DATOS PERSONALES DEL ESTUDIANTE
%========================================
\section{Datos Personales del Estudiante}

\begin{table}[h]
\centering
\begin{tabularx}{\textwidth}{|l|X|}
\hline
\textbf{Nombre y Apellidos} & Alejandro Vialard Santana \\
\hline
\textbf{DNI} & XXXXXXXXX \\
\hline
\textbf{Titulación} & Grado en Ingeniería Informática en Ingeniería del Software \\
\hline
\textbf{Correo electrónico} & avialards@alumnos.unex.es \\
\hline
\end{tabularx}
\end{table}

%========================================
% 2. DATOS DE LA EMPRESA/ENTIDAD
%========================================
\section{Datos de la Empresa/Entidad}

\subsection{Información General}

\begin{table}[h]
\centering
\begin{tabularx}{\textwidth}{|l|X|}
\hline
\textbf{Nombre de la Entidad} & Grupo de Investigación GICIS - Universidad de Extremadura \\
\hline
\textbf{Actividad} & Investigación en Ciberseguridad e Inteligencia Artificial aplicada a sistemas de seguridad informática \\
\hline
\textbf{Dirección} & Escuela Politécnica de Cáceres, Avda. de la Universidad s/n, 10003 Cáceres \\
\hline
\textbf{Departamento} & Grupo de Investigación en Ciberseguridad e Inteligencia de Sistemas (GICIS) \\
\hline
\end{tabularx}
\end{table}

\subsection{Datos del Tutor en la Entidad}

\begin{table}[h]
\centering
\begin{tabularx}{\textwidth}{|l|X|}
\hline
\textbf{Nombre} & [Nombre del Tutor] \\
\hline
\textbf{Cargo} & Investigador Principal / Profesor \\
\hline
\textbf{Correo electrónico} & [correo@unex.es] \\
\hline
\textbf{Teléfono de contacto} & [Teléfono] \\
\hline
\end{tabularx}
\end{table}

\subsection{Calendario y Características de las Prácticas}

\begin{table}[h]
\centering
\begin{tabularx}{\textwidth}{|l|X|}
\hline
\textbf{Fecha de inicio} & 24 de noviembre de 2025 \\
\hline
\textbf{Fecha de finalización} & 20 de enero de 2026 \\
\hline
\textbf{Número total de horas} & 150 horas \\
\hline
\textbf{Carácter de las prácticas} & No remuneradas \\
\hline
\end{tabularx}
\end{table}

\newpage

%========================================
% 3. ACTIVIDADES DESARROLLADAS
%========================================
\section{Actividades Desarrolladas}

\subsection{Descripción General del Proyecto}

Durante el periodo de prácticas en el Grupo de Investigación GICIS, he participado en un proyecto de investigación centrado en el diseño y desarrollo de un \textbf{agente autónomo basado en Modelos de Lenguaje de Gran Escala (LLM)} aplicado al ámbito de la \textbf{ciberseguridad ofensiva}.

El objetivo del proyecto es desarrollar un sistema capaz de realizar tareas de auditoría de seguridad de forma autónoma, utilizando la capacidad de razonamiento de los LLMs para planificar y ejecutar estrategias de pentesting. Este tipo de sistemas representa una línea de investigación emergente que combina las capacidades de razonamiento y generación de texto de los modelos de lenguaje con la ejecución de acciones en entornos reales.

La arquitectura del sistema se compone de tres capas principales:

\begin{itemize}
    \item \textbf{Capa de ejecución (Sandbox):} Un entorno aislado basado en contenedores Docker donde el agente puede ejecutar comandos de forma segura sin afectar al sistema host.
    
    \item \textbf{Capa de inteligencia (Backend):} Un servidor que gestiona la comunicación con múltiples proveedores de LLMs, implementa la lógica de Tool Calling y mantiene el estado de las conversaciones.
    
    \item \textbf{Capa de interacción (Frontend):} Una interfaz web que permite a los investigadores monitorizar el comportamiento del agente, ajustar parámetros y analizar los resultados de las auditorías.
\end{itemize}

\subsection{Actividades Realizadas por Semanas}

\subsubsection{Semanas 1-2 (24 noviembre -- 2 diciembre 2025): Infraestructura del Sandbox}

Durante las dos primeras semanas me centré en la construcción del entorno de ejecución aislado, que constituye la base técnica sobre la que el agente podrá operar de forma segura.

\paragraph{Diseño y construcción de la imagen Docker}

Se diseñó una imagen Docker personalizada basada en \textbf{Kali Linux}, seleccionada por su amplio catálogo de herramientas de auditoría preinstaladas. La imagen incluye:

\begin{itemize}
    \item Herramientas de reconocimiento: \texttt{nmap}, \texttt{masscan}, \texttt{nikto}.
    \item Herramientas de explotación: \texttt{metasploit}, \texttt{sqlmap}, \texttt{hydra}.
    \item Utilidades de red: \texttt{curl}, \texttt{wget}, \texttt{httpie}.
    \item Analizadores de APIs: \texttt{LinkFinder}, \texttt{nuclei}.
\end{itemize}

\paragraph{Desarrollo de la API de comunicación}

Se implementó una \textbf{API REST en Python utilizando Flask} que actúa como interfaz entre el sistema externo y el contenedor. Las funcionalidades principales incluyen:

\begin{itemize}
    \item Ejecución remota de comandos con captura de stdout y stderr.
    \item Gestión del directorio de trabajo actual (soporte para navegación con \texttt{cd}).
    \item Retorno estructurado de outputs incluyendo códigos de salida.
    \item Timeouts configurables para comandos de larga duración.
\end{itemize}

\paragraph{Orquestación del laboratorio de pruebas}

Mediante \textbf{Docker Compose} se desplegó una red aislada con segmento \texttt{10.5.0.0/24} que integra el agente con tres entornos vulnerables de referencia:

\begin{itemize}
    \item \textbf{DVWA (Damn Vulnerable Web Application):} Aplicación web diseñada para practicar vulnerabilidades clásicas como SQL Injection, XSS y CSRF.
    \item \textbf{Metasploitable 2:} Máquina virtual con múltiples servicios vulnerables a nivel de infraestructura (FTP, SSH, Samba, etc.).
    \item \textbf{OWASP Juice Shop:} Aplicación web moderna con vulnerabilidades representativas del OWASP Top 10.
\end{itemize}

\paragraph{Validación del entorno}

Se realizaron pruebas de comunicación y se ejecutó con éxito un ataque automatizado de \textbf{SQL Injection contra DVWA}, validando que el agente podía:
\begin{itemize}
    \item Realizar login automático en la aplicación.
    \item Evadir tokens CSRF mediante scripts personalizados.
    \item Extraer información de la base de datos mediante inyección SQL.
\end{itemize}

\subsubsection{Semana 3 (3 -- 4 diciembre 2025): Planificación y Coordinación}

Semana dedicada a la preparación de la siguiente fase del proyecto, con énfasis en la coordinación con recursos externos y la planificación técnica.

\begin{itemize}
    \item \textbf{Reunión con Computaex:} Se mantuvo una reunión con Daniel Flores para consultar la disponibilidad de GPUs en el supercomputador \textbf{Lusitania}. Se identificaron nodos con 2 NVIDIA RTX A4500 (20GB VRAM cada una) como objetivo para el despliegue de modelos de mayor escala.
    
    \item \textbf{Análisis de requisitos:} Se definieron las interfaces API necesarias, las estructuras de datos para el almacenamiento de conversaciones y los flujos de comunicación entre componentes.
    
    \item \textbf{Selección de modelos LLM:} Se elaboró un análisis de modelos compatibles con diferentes configuraciones de hardware, evaluando aspectos como ventana de contexto, soporte para Tool Calling y rendimiento en tareas de razonamiento.
    
    \item \textbf{Organización del repositorio:} Se estructuró el repositorio Git con ramas alineadas a las tareas planificadas, siguiendo un flujo de trabajo basado en feature branches.
\end{itemize}

\subsubsection{Semana 4 (5 -- 11 diciembre 2025)}

Periodo sin avances significativos en el proyecto debido a la priorización de entregas académicas de la asignatura de Procesamiento de Big Data (PBD).

\subsubsection{Semanas 5-6 (12 -- 23 diciembre 2025): Desarrollo del Backend e Integración con LLMs}

Periodo intensivo de desarrollo que abarcó tanto la investigación teórica como la implementación práctica del núcleo del sistema.

\paragraph{Investigación de APIs de LLMs}

Se realizó un estudio exhaustivo del formato de comunicación con modelos de lenguaje, incluyendo:

\begin{itemize}
    \item \textbf{Estructura de mensajes:} Análisis del sistema de roles (system, user, assistant) y gestión del contexto conversacional.
    \item \textbf{Tool Calling (Function Calling):} Estudio del mecanismo que permite a los LLMs invocar funciones externas, incluyendo la definición de esquemas JSON y el flujo de parámetros.
    \item \textbf{Diferencias entre proveedores:} Pruebas con OpenAI, Google Gemini, Anthropic, DeepSeek y OpenRouter para identificar particularidades en formato de respuestas y manejo de errores.
\end{itemize}

\paragraph{Implementación del Backend}

Se desarrolló un servidor \textbf{FastAPI} con las siguientes características:

\begin{itemize}
    \item \textbf{Endpoint de descubrimiento:} \texttt{/models/discover} que consulta múltiples proveedores en paralelo y devuelve una lista unificada de modelos disponibles.
    \item \textbf{Endpoint de chat:} \texttt{/chat} que implementa una interfaz universal compatible con tool calling, abstrayendo las diferencias entre proveedores.
    \item \textbf{Sistema de ranking:} Integración de puntuaciones ELO basadas en el LMArena Leaderboard, con fuzzy matching para identificar modelos por nombre.
    \item \textbf{Streaming de respuestas:} Soporte para Server-Sent Events (SSE) que permite visualizar las respuestas token a token.
\end{itemize}

\paragraph{Desarrollo del Frontend Web}

Se seleccionó un stack tecnológico moderno para la interfaz:

\begin{itemize}
    \item \textbf{Next.js 14:} Framework React con App Router para optimización y enrutamiento.
    \item \textbf{Tailwind CSS:} Framework de utilidades para estilos consistentes.
    \item \textbf{Shadcn/UI:} Biblioteca de componentes accesibles basada en Radix UI.
    \item \textbf{Zustand:} Gestión de estado global ligera y persistente.
\end{itemize}

Se implementó una interfaz de chat con renderizado de Markdown, resaltado de código y selección dinámica de modelos.

\paragraph{Diseño de la capa de persistencia}

Se diseñó una estructura de base de datos basada en \textbf{árboles N-arios} que permite:

\begin{itemize}
    \item Almacenar conversaciones de forma relacional.
    \item Implementar ``branching'' para explorar diferentes caminos de respuesta sin perder el historial original.
    \item Vincular el estado del sandbox a cada nodo del árbol de conversación.
\end{itemize}

\paragraph{Integración Sandbox-LLM}

Se desarrolló la conexión entre el modelo de lenguaje y el entorno de ejecución:

\begin{itemize}
    \item Cada conversación posee su propia carpeta aislada en el sandbox.
    \item Se implementó un sistema de versionado con Git que realiza commits automáticos por cada acción.
    \item Al navegar entre ramas del chat, el sandbox ejecuta \texttt{git checkout} al estado correspondiente.
\end{itemize}

\paragraph{Investigación de modelos locales}

Se analizó la viabilidad de ejecutar modelos en hardware limitado (RTX 3050, 4GB VRAM):

\begin{itemize}
    \item Pruebas con formatos de cuantización GGUF y AWQ (4-bit, 3-bit).
    \item Evaluación de modelos pequeños: Phi-3 Mini, Qwen-2.5-Coder, Gemma 2B.
    \item Benchmarking de velocidad de inferencia y consumo de memoria.
\end{itemize}

\paragraph{Acceso a infraestructura HPC}

Se estableció la primera conexión al supercomputador Lusitania de Computaex, estudiando su sistema de colas (SLURM) para planificar el despliegue futuro de modelos especializados en ciberseguridad.

\subsubsection{Semanas 7-8 (8 -- 20 enero 2026): Autonomía, Optimización y Documentación}

Fase final del proyecto centrada en dotar al agente de capacidades autónomas y en la documentación del trabajo realizado.

\paragraph{Gestión de procesos asíncronos}

Se implementaron funcionalidades para manejar procesos de larga duración:

\begin{itemize}
    \item Ejecución en segundo plano de comandos que no requieren respuesta inmediata.
    \item Panel de monitorización de procesos activos con visualización de salida en tiempo real.
    \item Sistema de callbacks que notifica automáticamente al LLM cuando un proceso finaliza.
    \item Funciones para detener procesos en ejecución desde la interfaz.
\end{itemize}

\paragraph{Optimización de la gestión de contexto}

Para mitigar las limitaciones de ventana de contexto de algunos modelos:

\begin{itemize}
    \item \textbf{Resumen automático:} Las salidas extensas de comandos son procesadas por un LLM auxiliar para generar resúmenes concisos.
    \item \textbf{Mecanismo de ``Zoom In'':} Permite al modelo expandir bajo demanda mensajes previamente resumidos.
    \item \textbf{Límite de contexto:} Sistema de recorte automático que mantiene los últimos 30 mensajes activos.
    \item \textbf{Memoria selectiva:} Preservación de mensajes marcados como relevantes más allá del límite general.
\end{itemize}

\paragraph{Sistema de planificación agéntica}

Se desarrolló un sistema de misiones que permite al agente mantener objetivos a largo plazo:

\begin{itemize}
    \item El LLM puede crear planes estructurados con subtareas y dependencias.
    \item Interfaz visual (Mission Board) con checkboxes para seguimiento del progreso.
    \item Memoria persistente independiente del contexto conversacional para almacenar hallazgos importantes.
    \item Endpoints REST (\texttt{/missions}) para crear, actualizar y consultar planes.
\end{itemize}

\paragraph{Optimización del comportamiento}

Se realizó una reingeniería del prompt del sistema para mejorar el comportamiento autónomo:

\begin{itemize}
    \item Instrucciones que fomentan la perseverancia ante fallos.
    \item Directrices para analizar errores y adaptar estrategias.
    \item Incentivos para buscar alternativas creativas cuando el enfoque directo no funciona.
\end{itemize}

\paragraph{Testing y documentación}

La fase final incluyó:

\begin{itemize}
    \item Ejecución de pruebas funcionales con escenarios complejos de reconocimiento y explotación.
    \item Generación de diagramas de arquitectura y flujos de datos.
    \item Redacción de documentación técnica del proyecto.
    \item Medición de métricas de rendimiento (tiempos de respuesta, tasas de éxito).
\end{itemize}

\newpage

\subsection{Tecnologías y Herramientas Utilizadas}

A lo largo del proyecto se han empleado diversas tecnologías que abarcan múltiples áreas de la ingeniería del software:

\subsubsection{Infraestructura y DevOps}

\begin{itemize}
    \item \textbf{Docker:} Contenedorización del entorno de ejecución y de los servicios del proyecto.
    \item \textbf{Docker Compose:} Orquestación del laboratorio de pruebas con múltiples contenedores.
    \item \textbf{Kali Linux:} Sistema operativo base del sandbox con herramientas de pentesting.
    \item \textbf{Git:} Control de versiones del código y del estado del sandbox.
\end{itemize}

\subsubsection{Desarrollo Backend}

\begin{itemize}
    \item \textbf{Python 3.11:} Lenguaje principal para el desarrollo del servidor y scripts.
    \item \textbf{FastAPI:} Framework web asíncrono de alto rendimiento.
    \item \textbf{SQLite:} Base de datos ligera para persistencia de conversaciones.
    \item \textbf{Server-Sent Events:} Protocolo para streaming de respuestas en tiempo real.
\end{itemize}

\subsubsection{Desarrollo Frontend}

\begin{itemize}
    \item \textbf{Next.js 14:} Framework React con renderizado del lado del servidor.
    \item \textbf{TypeScript:} Superset de JavaScript con tipado estático.
    \item \textbf{Tailwind CSS:} Framework de utilidades CSS para diseño responsivo.
    \item \textbf{Shadcn/UI:} Componentes de interfaz accesibles y personalizables.
    \item \textbf{Zustand:} Biblioteca de gestión de estado minimalista.
    \item \textbf{TanStack Query:} Gestión de datos asíncronos y caché.
\end{itemize}

\subsubsection{Inteligencia Artificial}

\begin{itemize}
    \item \textbf{APIs comerciales:} OpenAI (GPT-4), Google Gemini, Anthropic Claude, DeepSeek.
    \item \textbf{Agregadores:} OpenRouter para acceso unificado a múltiples proveedores.
    \item \textbf{Modelos locales:} Phi-3 Mini, Qwen-2.5-Coder, Gemma 2B en formatos cuantizados.
    \item \textbf{Tool Calling:} Mecanismo de invocación de funciones desde el LLM.
\end{itemize}

\subsubsection{Herramientas de Ciberseguridad}

\begin{itemize}
    \item \textbf{Nmap:} Escaneo de redes, descubrimiento de hosts y detección de servicios.
    \item \textbf{SQLMap:} Automatización de detección y explotación de SQL Injection.
    \item \textbf{Metasploit:} Framework de pruebas de penetración y desarrollo de exploits.
    \item \textbf{Hydra:} Herramienta de fuerza bruta para múltiples protocolos.
    \item \textbf{Nuclei:} Escáner de vulnerabilidades basado en plantillas YAML.
\end{itemize}

\newpage

\subsection{Valoración de las Tareas y Relación con la Titulación}

Las prácticas han permitido aplicar de forma integrada conocimientos adquiridos en múltiples asignaturas del Grado en Ingeniería Informática:

\begin{itemize}
    \item \textbf{Ingeniería del Software:} Aplicación de metodologías ágiles con sprints y entregas incrementales. Gestión del proyecto mediante issues y control de versiones con Git. Documentación técnica siguiendo estándares profesionales.
    
    \item \textbf{Arquitectura de Software:} Diseño de un sistema distribuido con separación clara de responsabilidades entre capas (sandbox, backend, frontend). Comunicación mediante APIs REST y protocolos de streaming.
    
    \item \textbf{Bases de Datos:} Diseño de esquemas no convencionales utilizando estructuras de árbol N-ario para el almacenamiento de conversaciones con capacidad de ramificación.
    
    \item \textbf{Redes de Computadores:} Configuración de redes virtuales aisladas, asignación de direcciones IP y comunicación entre contenedores Docker.
    
    \item \textbf{Seguridad Informática:} Aplicación práctica de conceptos de pentesting, identificación de vulnerabilidades web (OWASP Top 10) y uso de herramientas de auditoría.
    
    \item \textbf{Inteligencia Artificial:} Integración de modelos de lenguaje, diseño de prompts efectivos y desarrollo de sistemas agénticos con capacidad de planificación.
\end{itemize}

\subsection{Aprendizaje Obtenido}

Durante el periodo de prácticas he adquirido conocimientos y habilidades en las siguientes áreas:

\subsubsection{Arquitectura de Sistemas con Inteligencia Artificial}

\begin{itemize}
    \item Comprensión profunda de la estructura de comunicación con LLMs: roles, tokens, ventanas de contexto.
    \item Dominio del mecanismo de Tool Calling y su implementación en diferentes proveedores.
    \item Estrategias para mitigar limitaciones de los modelos: alucinaciones, pérdida de contexto, falta de persistencia.
    \item Importancia del diseño de prompts (prompt engineering) en el comportamiento del sistema.
\end{itemize}

\subsubsection{Desarrollo Full-Stack Moderno}

\begin{itemize}
    \item Implementación de backends asíncronos de alto rendimiento con FastAPI.
    \item Desarrollo de frontends reactivos con Next.js y gestión de estado compleja.
    \item Comunicación en tiempo real mediante Server-Sent Events.
    \item Patrones de diseño para aplicaciones web escalables.
\end{itemize}

\subsubsection{Ciberseguridad Aplicada}

\begin{itemize}
    \item Metodologías de pentesting: reconocimiento, enumeración, explotación.
    \item Uso práctico de herramientas especializadas (nmap, sqlmap, metasploit).
    \item Configuración de entornos de laboratorio seguros y aislados.
    \item Comprensión de vulnerabilidades web comunes y sus vectores de explotación.
\end{itemize}

\subsubsection{Gestión de Proyectos de Investigación}

\begin{itemize}
    \item Coordinación con recursos compartidos como infraestructuras HPC.
    \item Documentación académica y técnica con rigor metodológico.
    \item Planificación a largo plazo con adaptación ante imprevistos.
    \item Investigación bibliográfica y análisis del estado del arte.
\end{itemize}

\newpage

\subsection{Valoración General y Conclusiones}

\subsubsection{Aspectos Positivos}

Esta experiencia de prácticas ha sido \textbf{extraordinariamente enriquecedora} por múltiples razones:

\begin{itemize}
    \item \textbf{Proyecto innovador:} Trabajar en la intersección de IA y ciberseguridad representa una oportunidad única de contribuir a un campo emergente con gran potencial de impacto tanto en investigación como en aplicaciones prácticas.
    
    \item \textbf{Autonomía técnica:} He tenido libertad para tomar decisiones de diseño y arquitectura, lo que ha fomentado un aprendizaje profundo y responsable.
    
    \item \textbf{Stack tecnológico actual:} Las tecnologías utilizadas (Next.js 14, FastAPI, Docker, LLMs) representan el estado del arte en desarrollo de software, proporcionando competencias directamente aplicables en el mercado laboral.
    
    \item \textbf{Entorno de investigación:} Trabajar en un grupo de investigación proporciona una perspectiva diferente al desarrollo en empresa, con énfasis en la innovación, el rigor metodológico y la documentación.
\end{itemize}

\subsubsection{Retos Superados}

Durante las prácticas me he enfrentado a diversos desafíos técnicos:

\begin{itemize}
    \item \textbf{Integración multi-proveedor:} Cada API de LLM tiene particularidades; diseñar una capa de abstracción efectiva requirió un estudio profundo y pruebas exhaustivas.
    
    \item \textbf{Sincronización de estado:} Mantener coherencia entre el estado del chat y el sistema de archivos del sandbox mediante Git fue un reto de ingeniería complejo.
    
    \item \textbf{Optimización de contexto:} Desarrollar estrategias para que el modelo mantenga coherencia con ventanas de contexto limitadas requirió experimentación iterativa.
    
    \item \textbf{Comportamiento del agente:} Lograr que el LLM sea persistente y proactivo ante obstáculos requirió múltiples iteraciones de ingeniería de prompts.
\end{itemize}

\subsubsection{Conclusión}

En definitiva, las prácticas me han permitido consolidar conocimientos adquiridos durante la carrera, desarrollar competencias en tecnologías de vanguardia, cultivar el pensamiento crítico necesario para la toma de decisiones técnicas fundamentadas, y confirmar mi interés en la ciberseguridad e inteligencia artificial como área de especialización profesional.

Considero que esta experiencia ha sido altamente formativa y recomendable, proporcionándome herramientas, conocimientos y una visión práctica que complementan de forma significativa mi formación académica.

\newpage

%========================================
% REFERENCIAS
%========================================
\section{Referencias Utilizadas}

\begin{enumerate}
    \item \textbf{Documentación oficial de FastAPI:} \url{https://fastapi.tiangolo.com/}
    
    \item \textbf{Documentación de Next.js:} \url{https://nextjs.org/docs}
    
    \item \textbf{Shadcn/UI Components:} \url{https://ui.shadcn.com/}
    
    \item \textbf{Docker Documentation:} \url{https://docs.docker.com/}
    
    \item \textbf{OpenAI API Reference:} \url{https://platform.openai.com/docs/api-reference}
    
    \item \textbf{Google Gemini API:} \url{https://ai.google.dev/docs}
    
    \item \textbf{Anthropic Claude API:} \url{https://docs.anthropic.com/}
    
    \item \textbf{Kali Linux Tools:} \url{https://www.kali.org/tools/}
    
    \item \textbf{OWASP Testing Guide:} \url{https://owasp.org/www-project-web-security-testing-guide/}
    
    \item \textbf{LMArena Leaderboard:} \url{https://lmarena.ai/}
    
    \item \textbf{Zustand State Management:} \url{https://zustand-demo.pmnd.rs/}
    
    \item \textbf{TanStack Query:} \url{https://tanstack.com/query/latest}
    
    \item \textbf{DVWA Documentation:} \url{https://github.com/digininja/DVWA}
    
    \item \textbf{Metasploitable 2:} \url{https://sourceforge.net/projects/metasploitable/}
    
    \item \textbf{OWASP Juice Shop:} \url{https://owasp.org/www-project-juice-shop/}
\end{enumerate}

\newpage

%========================================
% 4. FECHA Y FIRMA
%========================================
\section{Fecha de Entrega y Firma}

\vspace{2cm}

\begin{center}
    \textbf{Fecha de entrega:} 27 de enero de 2026
\end{center}

\vspace{3cm}

\begin{center}
    \textbf{Firma del estudiante:}
\end{center}

\vspace{3cm}

\begin{center}
    \rule{8cm}{0.4pt}\\
    \vspace{0.5cm}
    Alejandro Vialard Santana
\end{center}

\end{document}
